\chapter{API Design and Implementation}

\section{API Design Strategy}

The API design separates actor responsibilities instead of forcing every actor
through one shared access model. This decision keeps merchant integration,
internal operations, and merchant dashboard access aligned with their real
security and data-visibility needs.

Detailed request and response contracts are maintained in a separate API
reference. This chapter focuses on the boundary design decisions that shape the
overall interface.

Three interface rules drive the design:

\begin{itemize}
  \item each actor group receives an access surface that matches its real
  operational role rather than being forced through one generic API;
  \item state-changing integration flows remain separate from human dashboard
  visibility and internal control flows;
  \item delayed callbacks and outbound notifications are treated as durable
  event exchanges rather than incidental side effects.
\end{itemize}

\section{API Families And Access Boundaries}

The family map below makes that split explicit by pairing each surface with its
primary actor, its trust model, and its design purpose.

\begin{longtable}{L{0.17\textwidth}L{0.19\textwidth}L{0.20\textwidth}L{0.34\textwidth}}
\caption{API families and access boundaries}\\
\toprule
\textbf{API family} & \textbf{Primary actor} & \textbf{Authentication model} & \textbf{Design purpose} \\
\midrule
\endfirsthead
\toprule
\textbf{API family} & \textbf{Primary actor} & \textbf{Authentication model} & \textbf{Design purpose} \\
\midrule
\endhead
Internal authentication and administration & Admin or Ops user & Internal session authentication & Supports internal sign-in, account management, and operational access control. \\
\bodyrowrule
Merchant transaction APIs & Merchant backend & Signed merchant integration credentials & Supports payment creation, payment lookup, refund creation, and refund lookup for merchant systems. \\
\bodyrowrule
Provider callback APIs & Provider callback source & Provider id, timestamp, and HMAC signature headers & Accepts external outcome evidence for payments and refunds. \\
\bodyrowrule
Internal operations APIs & Admin or Ops user & Internal session authentication & Supports merchant management, QR receiving account management, evidence review, reconciliation, notification recovery, and dashboard reads. \\
\bodyrowrule
Merchant dashboard APIs & Merchant dashboard user & Merchant session authentication & Supports merchant-scoped overview, analytics, record exploration, profile views, and credential metadata views. \\
\bodyrowrule
Outbound notification contract & Merchant callback endpoint & Signed gateway-to-merchant notification model & Delivers final payment and refund events to merchant-owned callback destinations. \\
\bottomrule
\end{longtable}

\section{Authentication Model Split}

As established by the system architecture, the interface preserves three access
boundaries: merchant integration requests, internal operational sessions, and
merchant-facing dashboard sessions. This split prevents two common design
problems: using human dashboard login for machine integration, and using
merchant integration credentials for interactive dashboard access.

\section{VietQR Payment API Contract}

Payment creation keeps the merchant request shape intentionally small, but the
response now carries the information needed for a merchant checkout page to show
a scannable VietQR code. A successful response includes:

\begin{itemize}
  \item \code{qr_reference}: a short gateway-generated transfer reference used
  for matching and support;
  \item \code{qr_content}: the generated VietQR payload that starts with the
  EMV merchant-presented QR prefix;
  \item \code{qr_image_base64}: a PNG data URL that can be rendered directly by
  a merchant checkout or dashboard detail view.
\end{itemize}

The gateway only returns these values after the merchant is active, has active
integration credentials, and has an active Ops-managed VietQR receiving account.
The API rejects non-VND or fractional pilot payments before QR generation so
that invalid checkout data is not persisted.

\section{Asynchronous API And Event Design}

The interface design must safely accept delayed external evidence and delayed
merchant notification delivery. For that reason, provider callbacks and
merchant notifications are treated as durable event flows rather than as simple
status updates.

When external evidence arrives, the system first preserves the evidence and then
applies state rules. If the evidence conflicts with the current internal state,
the system creates a reconciliation path rather than forcing an unsafe
automatic correction.

Provider callbacks are protected by \code{X-Provider-Id},
\code{X-Provider-Timestamp}, and \code{X-Provider-Signature}. The signature
covers the timestamp and raw request body, which prevents unauthenticated or
stale callback messages from changing payment or refund state.

Merchant notification delivery is also decoupled from transaction finality. A
payment or refund can be final even if merchant notification delivery fails
temporarily. This makes retry and manual recovery possible without rewriting the
underlying money-flow outcome.

\section{Error Model And Data Exposure Rules}

The API design uses stable business error codes so that calling systems and
internal users can distinguish validation failure, readiness failure,
authentication failure, and business-rule conflict.

Sensitive data exposure is minimized throughout the interface design:

\begin{itemize}
  \item raw merchant secrets are not exposed in ordinary read views;
  \item generated one-time passwords are shown only at the moment of creation
  or reset;
  \item merchant dashboard reads remain scoped to the signed-in merchant;
  \item QR receiving account management remains on the internal Ops surface,
  not the merchant dashboard surface;
  \item operational evidence is preserved for investigation without exposing
  secret material unnecessarily.
\end{itemize}
